\begin{claim}{D.33}{Proved}{Finite punctured-torus actions are detected by their primary homology}
Put $F=\widehat F(x,y)$ and $c=[x,y]$, and let
$T=\Out_{\rm lac}(F;\overline{\langle c\rangle})$ be the full
once-punctured-torus loop-outer group.  For every prime $p$, the
homomorphism induced by the named abelianization,
\[
              A_p:T\longrightarrow\GL_2(\Z_p),            \tag{D.33a}
\]
is injective on every finite $p$-subgroup of $T$.
The full abelianization homomorphism
$T\to\GL_2(\widehat\Z)$ is therefore injective on every finite
subgroup.  Every such subgroup has order dividing $24$.

The fixed-subgroup calculation underlying these assertions is also
explicit.  For an automorphism $\alpha$ of $\widehat F_r$ with
$\alpha^p=1$ whose induced automorphism on the maximal pro-$p$
quotient is the identity, put $E=\widehat F_r\rtimes_\alpha C_p$.
There is one kernel-conjugacy class of sections of $E\to C_p$,
and for its fixed subgroup $C$ the actual restriction map is an
isomorphism
\[
 H^1(\widehat F_r,\F_p)\ \xrightarrow{\ \sim\ }H^1(C,\F_p).
                                                        \tag{D.33b}
\]
For the punctured-torus loop structure and $r=2$, its cusp forces
$C=F$ and hence $\alpha=1$.
\end{claim}
\begin{proof}
We first prove (D.33b), keeping the multiplicative information in
the cohomology comparison.  Write $P=\widehat F_r^{(p)}$ and
$S=\{1\}\times C_p\leq E$.  The quotient homomorphism is
\[
 \pi:E\longrightarrow P\times C_p,\qquad
              (h,s^i)\longmapsto(\bar h,s^i).              \tag{D.33c}
\]
Restriction along $\widehat F_r\to P$ is an isomorphism on
mod-$p$ cohomology: in degree one both sides are the same
$r$-dimensional character group, and higher cohomology vanishes
because both groups have $p$-cohomological dimension one.
The Hochschild--Serre comparison for (D.33c) is consequently an
isomorphism on its second pages and gives a ring isomorphism
\[
 \pi^*:H^*(P\times C_p,\F_p)\xrightarrow{\sim}H^*(E,\F_p).
                                                        \tag{D.33d}
\]
All maps here are induced by the specified quotient, not by an
identification of the two groups.  In particular
$\dim H^n(E,\F_p)=r+1$ for every $n>1$.

Let $u\in H^2(C_p,\F_p)$ be the nonzero periodic class, and
inflate it to either side of (D.33d).  The K\"unneth calculation
for the direct product gives, for every $m\geq1$, an injective map
\[
 H^1(P,\F_p)\longrightarrow H^{2m+1}(E,\F_p),
                         \qquad\chi\longmapsto\pi^*(\chi u^m).
                                                        \tag{D.33e}
\]
For a subgroup $Q\leq E$ of order $p$, its projection to
$C_p$ is an isomorphism, and
$N_E(Q)=C_{\widehat F_r}(Q)\times Q$, exactly as in D.28.
The first coordinate of $\pi(Q)$ is trivial because $P$ is
torsion-free.  Thus every character $\pi^*\chi$ in (D.33e)
vanishes on $Q$.  Its restriction to this normalizer lies
entirely in $H^1(C_{\widehat F_r}(Q),\F_p)$, and multiplication
by $u^m$ places it in
\[
 H^1(C_{\widehat F_r}(Q),\F_p)\otimes H^{2m}(Q,\F_p).
                                                        \tag{D.33f}
\]

Apply the profinite Brown restriction theorem in the precise
form verified in \cite[Theorem 2 and \S1.1]{A-KochloukovaZalesskii2007}.
There is no $C_p\times C_p$ in $E$, since a finite subgroup
injects into $E/\widehat F_r=C_p$.  Its open subgroup
$\widehat F_r$ has $p$-cohomological dimension one.  The
finite-dimensional source makes the theorem's dense restriction
embedding an isomorphism onto a finite product.  If $k$ is
the number of conjugacy classes and
$d_Q=\dim H^1(C_{\widehat F_r}(Q),\F_p)$, its dimension formula is
\[
                   k+\sum_{[Q]}d_Q=r+1.                  \tag{D.33g}
\]
The injection (D.33e), followed by this same restriction map,
lands in the summands (D.33f).  Therefore
$\sum d_Q\geq r$.  Since $k\geq1$, (D.33g) forces
$k=1$ and $d_Q=r$.  For the specified $S$, (D.33e)--(D.33f)
then say exactly that restriction
$H^1(\widehat F_r,\F_p)\to H^1(C_{\widehat F_r}(S),\F_p)$
is injective; both spaces have dimension $r$, proving (D.33b).
As in D.28, conjugacy under $E$ equals conjugacy under the kernel,
because $E=\widehat F_r Q$ for any section $Q$.

Now take $r=2$ and assume that $\alpha$ preserves the conjugacy
class of $L=\overline{\langle c\rangle}$, where $c=[x,y]$.
Write its peripheral multiplier as $\lambda$.  In a class-two
Heisenberg quotient, the image of the commutator under an
abelianization matrix $A$ is the $\det(A)$ power of the
commutator.  The finite Heisenberg quotients for all moduli
therefore give the exact identity
\[
                        \lambda=\det(\alpha^{\rm ab}).   \tag{D.33h}
\]
In the present calculation $\alpha$ is the identity on $P$,
so $\lambda_p=1$.

The cusp normalizer has the exact sequence
\[
 1\longrightarrow L\longrightarrow N_E(L)
                        \longrightarrow C_p\longrightarrow1.
                                                        \tag{D.33i}
\]
The kernel is $L$ by A.7.  For a lift $a$ of the chosen
generator, write
$aca^{-1}=c^\lambda$ and $a^p=c^v$, so $(\lambda-1)v=0$.
The first coordinate of $\pi(a)$ centralizes the image $\bar c$
in $P$ and belongs to $\overline{\langle\bar c\rangle}$,
by the pro-$p$ peripheral centralizer assertion of A.7.
Thus it is $\bar c^w$ and $v_p=pw$.
This uses the whole cusp subgroup.  Its primitive pro-$p$
parameter can also be checked directly: in the mod-$p$ Magnus
expansion, $[x,y]-1$ starts with $XY-YX$.  A $p$th power
starts in degree divisible by $p$.  At $p=2$, the possible
degree-two square is $(aX+bY)^2$; its $X^2$ and $Y^2$
coefficients prevent it from equaling $XY+YX$.  Hence the
commutator is not a $p$th power.

Since $\alpha^p=1$, its multiplier satisfies $\lambda^p=1$.
Set $N_\lambda=1+\lambda+\cdots+\lambda^{p-1}$.
At a prime $q$ with $\lambda_q\ne1$, both
$(N_\lambda)_q$ and $v_q$ are zero.  Where $\lambda_q=1$
the norm is $p$; $v_q$ is divisible by $p$, using $v_p=pw$
at $q=p$ and invertibility of $p$ at all other primes.
Consequently there is $z\in\widehat\Z$ with
$N_\lambda z=-v$.  The element $ac^z$ has order $p$
and centralizes $c^{e_p\Z_p}$.
The unique section-conjugacy class just proved places a
conjugate $g c^{e_p\Z_p}g^{-1}$ inside
$C=C_F(S)$.

Every character of $C$ vanishes on $g c^{e_p}g^{-1}$:
by (D.33b) it is the restriction of a character of $F$,
and $c$ is a commutator.  We prove that this forces $C=F$.
Its mod-$p$ first-cohomology dimension is two.  If $C$ is
open, Schreier gives $1+[F:C]=2$.  If $C$ is proper and
not open, the restriction-colimit proof of D.32 shows that
its index has infinite $p$-part.  Choose an open $U\geq C$
with $p\mid n=[F:U]$, and take the degree-$n$ covering of
the punctured torus corresponding to $U\cap F(x,y)$.
The chosen cusp has a boundary lift of degree $d$ prime
to $p$, since its subgroup contains the full chosen
$c^{e_p\Z_p}$.  The comparison with a discrete boundary
generator is exactly (D.32d), with a dense discrete
representative for the profinite conjugator $g$.
There must be at least two boundary lifts: a unique lift
would have degree $n$, divisible by $p$.
The selected boundary has nonzero homology in this cover,
since the sole relation among the at least two boundary
classes is their sum.  A character of $U$ detecting this
boundary restricts to a character of $C$ detecting
$g c^{e_p}g^{-1}$; the value is multiplied by
$d^{-1}\bmod p$, as computed in D.32.  This is the required
contradiction.  Hence $C=F$ and $\alpha=1$.

To apply this to (D.33a), let an outer element have order
$p$ and $A_p=I$.  D.27 supplies an actual order-$p$
automorphism representing it.  Its induced pro-$p$
automorphism has identity abelianization.  The Magnus
argument in D.10 and D.25 makes this finite-order
automorphism the identity: a first nonzero homogeneous
deviation in the $\Z_p$ Magnus algebra would be multiplied
by $p$ under the $p$th iterate, and the homogeneous module
has no $p$-torsion.  Thus (D.33c) applies to this lift
directly; there is no subsequent conjugation which could
change its order.  The cusp argument proves that the lift,
and hence its outer class, is the identity.
The kernel of (D.33a) on a finite $p$-subgroup cannot
contain an order-$p$ element, so it is trivial.  Likewise,
the kernel of full abelianization on any finite subgroup
would contain an element of some prime order and is trivial.

We finish with the numerical bound on these finite groups.
For $p>3$, a matrix of order $p$ over $\Q_p$ has a primitive
$p$th root among its eigenvalues.  The polynomial
$\Phi_p(t)$ is irreducible of degree $p-1$ over $\Q_p$,
as seen by Eisenstein after $t\mapsto t+1$.  It cannot be
the factor of a degree-two characteristic polynomial.
Thus no such prime divides the order of a finite subgroup
of $T$.  A finite $3$-subgroup injects by $A_3$ into
$\GL_2(\Z_3)$ and then injects modulo $3$: the kernel
$1+3M_2(\Z_3)$ is torsion-free by the matrix logarithm.
Since $|\GL_2(\F_3)|=48$, its order divides $3$.

For completeness, a finite $2$-subgroup $K$ of
$\GL_2(\Q_2)$ has order at most eight.  The same
Eisenstein argument for $\Phi_{2^a}$ excludes order
$2^a$ for $a\geq3$.  Every matrix of order four has
minimal polynomial $t^2+1$, so its square is $-I$.
Consequently the projective image of $K$ has exponent
two.  Extend the matrices to the algebraic closure
$\overline{\Q}_2$; the scalar kernel is unchanged.
An exponent-two subgroup of $\PGL_2(\overline{\Q}_2)$
has order at most four.  To prove this
last bound, conjugate a nonidentity element to
$z\mapsto-z$.  Every commuting transformation preserves
its pair of fixed points, so is $z\mapsto az$ or
$z\mapsto a/z$.  The scaling subgroup has at most two
elements under the exponent-two condition and has index
at most two.  The scalar kernel of $K$ is contained in
$\{I,-I\}$, the finite $2$-torsion of $\Q_2^\times$.
Thus $|K|\leq8$.  Apply this to a Sylow $2$-subgroup
through $A_2$.  The Sylow bounds and the exclusion of
other primes give $|G|\mid 8\cdot3=24$.

The program \texttt{code/part\_c/torus\_homology\_d33.py}
checks the Magnus degree-two coefficients, the periodic
cohomology multiplication used in (D.33e), and the
one-cusp restriction argument on every connected finite
cover through degree six.  The matrices and boundary
characters are obtained from actual Schreier walks.
The bound and comparison maps above are proved for the
full profinite groups, independently of the bounded checks.
\end{proof}
